\uuid{MFFm}
\exo7id{6921}
\titre{exo7 6921}
\auteur{ruette}
\organisation{exo7}
\datecreate{2013-01-24}
\isIndication{false}
\isCorrection{false}
\chapitre{Probabilité continue}
\sousChapitre{Loi normale}
\cpge{bcpst2.variables-aleatoires-densite.c, bcpst2.variables-aleatoires-densite.d, bcpst2.variables-aleatoires-densite.e}
\module{Probabilité et statistique}
\niveau{L2}
\difficulte{}

\contenu{
\begin{enumerate}
    \item \question{Soit $X$ une variable aléatoire suivant une loi normale $N(m,\sigma^2)$.
Quelle est la probabilité que $X$ soit supérieure à $m$ ?}
    \item \question{Soit  $Y$ une variable aléatoire suivant une loi normale $N(0,1)$.
Quelle est la loi de $-Y$ ? la loi de $\sigma Y+m$ ?}
    \item \question{En utilisant la table de $\mathcal{N}(0,1)$, déterminer 
$P(0\le Y\le 0,8)$,
$P(-0,6\le Y\le 0)$ et $P(Y\le 0,8)$.}
\end{enumerate}
}
